% GENERATED by scripts/19_build_submission_paper.py; edit the canonical theory source.
\section{Detailed proofs}
\label{rev:sec-detailed-proofs}

This appendix expands the measure-theoretic steps that are easy to hide inside
phrases such as ``interchange the limit and the conditional expectation.''
It covers the general results used repeatedly in the main text.  The
model-specific OU, CIR, and special-function calculations remain next to their
statements because they are direct substitutions into the displayed
differential equations.

\subsection{Copula kernels, density versions, and partial derivatives}
\label{rev:app-copula-kernel-proof}

\begin{proof}[Detailed proof of the kernel and derivative claims in
Proposition~\ref{rev:prop-density-copula-consistency}]
Fix $s<t$ and suppress the time subscripts temporarily.  For the selected
jointly Borel density $c$, define
\[
K(u,B)=\int_Bc(u,v)\dd v,
\qquad
H^c(u,y)=K(u,(0,y])=\int_0^yc(u,v)\dd v.
\]
Row normalization makes $B\mapsto K(u,B)$ a probability measure for every
$u$.  Joint measurability of $c$, first for intervals and then by the monotone
class theorem for all $B\in\mathcal B(I)$, makes $u\mapsto K(u,B)$ Borel.
Column normalization and Tonelli's theorem give, for every Borel $B$,
\begin{align}
\int_IK(u,B)\dd u
&=\int_B\left\{\int_Ic(u,v)\dd u\right\}\dd v
=\Leb(B),
\label{rev:app-kernel-invariance}
\end{align}
so the selected kernel preserves the uniform law.

For every $x,y\in[0,1]$, the copula reconstructed from $c$ satisfies
\begin{equation}
C(x,y)=\int_0^xH^c(a,y)\dd a.
\label{rev:app-copula-kernel-reconstruction}
\end{equation}
For fixed $y$, the function $H^c(\cdot,y)$ is measurable and bounded by one.
The fundamental theorem for Lebesgue integrals therefore implies
\begin{equation}
\partial_1C(u,y)=H^c(u,y)
\quad\text{for $\Leb$-a.e. }u.
\label{rev:app-partial-from-kernel}
\end{equation}
This is an equality for each fixed $y$; its exceptional set may initially
depend on $y$.  Because $H^c(u,\cdot)$ and every conditional-CDF version are
right-continuous, equality at all rational $y$ outside one null set extends to
all $y\in[0,1]$ outside that same source-state null set.

At any source state $u$ where $a\mapsto H^c(a,y)$ is continuous, the pointwise
upgrade follows directly from the integral difference quotient:
\begin{align}
\frac{C(u+h,y)-C(u,y)}{h}-H^c(u,y)
&=\frac1h\int_u^{u+h}\{H^c(a,y)-H^c(u,y)\}\dd a,\notag\\
\left|\frac{C(u+h,y)-C(u,y)}{h}-H^c(u,y)\right|
&\le \sup_{|a-u|\le |h|}|H^c(a,y)-H^c(u,y)|
\longrightarrow0.
\label{rev:app-pointwise-partial-upgrade}
\end{align}
Thus a continuity condition on the integrated density row, not merely existence
of a $\Leb^2$-a.e. density, is what upgrades
\eqref{rev:app-partial-from-kernel} to a pointwise derivative identity.  The
modified partial Dini construction of \citet{FangJiangLiuYang2020} supplies
another pointwise CDF-valued version $\mathfrak D_1C(u,\cdot)$ satisfying the
same reconstruction formula.  Applying
\eqref{rev:app-partial-from-kernel} to both versions shows that
$\mathfrak D_1C(u,\cdot)=H^c(u,\cdot)$ for a.e. $u$, although equality can fail
on a source-null set.  These uses of absolute continuity and differentiation
are instances of the results in \citet[Chapter~3]{Folland1999}.

Now fix $0\le s<r<t$ and $u\in I$.  Since
\eqref{rev:copula-density-CK} holds for a.e. $v$, equality of its integrals over
an arbitrary $B\in\mathcal B(I)$ and Tonelli's theorem give
\begin{align}
K_{s,t}(u,B)
&=\int_B\int_Ic_{s,r}(u,w)c_{r,t}(w,v)\dd w\dd v\notag\\
&=\int_Ic_{s,r}(u,w)
       \left\{\int_Bc_{r,t}(w,v)\dd v\right\}\dd w\notag\\
&=\int_IK_{r,t}(w,B)K_{s,r}(u,\dd w).
\label{rev:app-density-to-kernel-CK}
\end{align}
This proves the kernel Chapman--Kolmogorov equation for every selected source
$u$ and every Borel $B$.  Conversely, if both sides are absolutely continuous
in the target variable, equality for every $B$ identifies their
Radon--Nikodym densities only for a.e. $v$.

Finally, let $C=C_{s,r}$ and $D=C_{r,t}$.  The functions
\[
w\longmapsto\int_0^xc_{s,r}(a,w)\dd a,
\qquad
w\longmapsto\int_0^yc_{r,t}(w,b)\dd b
\]
are measurable a.e. representatives of $\partial_2C(x,w)$ and
$\partial_1D(w,y)$, respectively.  Substitution into
\eqref{rev:star-product-density-version}, followed by Tonelli and density CK,
gives
\begin{align}
(C_{s,r}*C_{r,t})(x,y)
&=\int_I\int_0^xc_{s,r}(a,w)\dd a
             \int_0^yc_{r,t}(w,b)\dd b\dd w\notag\\
&=\int_0^x\int_0^y
     \left\{\int_Ic_{s,r}(a,w)c_{r,t}(w,b)\dd w\right\}\dd b\dd a\notag\\
&=\int_0^x\int_0^yc_{s,t}(a,b)\dd b\dd a
=C_{s,t}(x,y).
\label{rev:app-star-product-integral}
\end{align}
All changes in integration order are justified by nonnegativity.  Conversely,
equality of these copula CDFs identifies their two-dimensional densities only
$\Leb^2$-a.e. and hence, by Fubini, yields conditional density composition only
for a.e. source state.  This proves why star-product consistency alone is weaker
than the selected pointwise kernel condition in
Assumption~\ref{rev:ass-copula-density}.
\end{proof}

\subsection{Short-time kernels acting on moving functions}

\begin{lemma}[Moving-function kernel lemma]
\label{rev:app-kernel-lemma}
Fix $t\ge0$ and an interior state $u\in I$.  Write
\[
K_h(u,\dd v)=c_{t,t+h}(u,v)\dd v,
\qquad h>0,
\]
so that $K_h(u,I)=1$.  The following statements hold.
\begin{enumerate}[label=(\alph*)]
\item Suppose that
\begin{equation}
K_h\bigl(u,\{v:|v-u|\ge\delta\}\bigr)\longrightarrow0
\quad\text{for every }\delta>0.
\label{rev:kernel-stochastic-continuity}
\end{equation}
If $g$ is bounded and continuous at $u$, then
\begin{equation}
\int_Ig(v)K_h(u,\dd v)\longrightarrow g(u).
\label{rev:kernel-fixed-function-limit}
\end{equation}
\item In addition to part~(a), let $r_h:I\to\R$ be measurable and suppose
$g$ is bounded and continuous at $u$ and
\begin{equation}
\lVert r_h-g\rVert_\infty\longrightarrow0.
\label{rev:moving-uniform-limit}
\end{equation}
Then
\begin{equation}
\int_Ir_h(v)K_h(u,\dd v)\longrightarrow g(u).
\label{rev:moving-kernel-limit}
\end{equation}
\item For an unbounded $g$, let $w:I\to[1,\infty)$ be measurable and define
\[
\lVert f\rVert_w=\sup_{v\in I}\frac{|f(v)|}{w(v)}.
\]
If
\begin{align}
\lVert r_h-g\rVert_w&\longrightarrow0,
\label{rev:moving-weighted-limit}\\
\sup_{0<h<h_0}\int_Iw(v)K_h(u,\dd v)&<\infty,
\label{rev:kernel-weight-bound}\\
\int_Ig(v)K_h(u,\dd v)&\longrightarrow g(u),
\label{rev:weighted-strong-continuity}
\end{align}
then \eqref{rev:moving-kernel-limit} still holds.
\end{enumerate}
\end{lemma}

\begin{proof}
For part~(a), split the integral at distance $\delta$ from $u$:
\begin{align*}
\left|\int_Ig(v)K_h(u,\dd v)-g(u)\right|
&\le \int_I|g(v)-g(u)|K_h(u,\dd v)\\
&\le \sup_{|v-u|<\delta}|g(v)-g(u)|
 +2\lVert g\rVert_\infty
 K_h\bigl(u,\{|v-u|\ge\delta\}\bigr).
\end{align*}
First let $h\downarrow0$ and use
\eqref{rev:kernel-stochastic-continuity}; then let $\delta\downarrow0$ and
use continuity of $g$ at $u$.  This proves
\eqref{rev:kernel-fixed-function-limit}, or equivalently
$K_h(u,\cdot)\Rightarrow\delta_u$.

For part~(b), row normalization gives
\begin{align*}
\left|\int_Ir_h(v)K_h(u,\dd v)-g(u)\right|
&\le \int_I|r_h(v)-g(v)|K_h(u,\dd v)\\
&\quad+\left|\int_Ig(v)K_h(u,\dd v)-g(u)\right|\\
&\le \lVert r_h-g\rVert_\infty
 +\left|\int_Ig(v)K_h(u,\dd v)-g(u)\right|.
\end{align*}
Both terms tend to zero by \eqref{rev:moving-uniform-limit} and part~(a).

For part~(c), replace the first term in the preceding bound by
\[
\int_I|r_h(v)-g(v)|K_h(u,\dd v)
\le \lVert r_h-g\rVert_w\int_Iw(v)K_h(u,\dd v).
\]
Equations \eqref{rev:moving-weighted-limit}--
\eqref{rev:weighted-strong-continuity} prove the claim.  This weighted version
is the relevant one for Normal, Student, or other quantiles that diverge at the
endpoints of $I$.
\end{proof}

Pointwise convergence $r_h(v)\to g(v)$ alone is not enough in this lemma: the
measure $K_h(u,\dd v)$ also changes with $h$.  Conditions
\eqref{rev:moving-uniform-limit} or
\eqref{rev:moving-weighted-limit}--\eqref{rev:kernel-weight-bound} provide the
uniform integrability needed for the moving integrand.

\subsection{Full integral proof of the quantile--generator drift identity}
\label{rev:app-drift-identity-proof}

\begin{proof}[Detailed proof of Proposition~\ref{rev:prop-generator-sign}]
Fix $t$ and put $q_t(v)=q(t,v)=Q_t(v)$ and
$\dot q_t(v)=\partial_tq(t,v)$.  Assume that all displayed first moments are
finite and that, for the state $u$ under consideration,
\begin{equation}
\mathcal A_tq_t(u)
=\lim_{h\downarrow0}
\frac{\int_Iq_t(v)c_{t,t+h}(u,v)\dd v-q_t(u)}{h}
\label{rev:app-generator-integral}
\end{equation}
exists.  Also assume either part~(b) or part~(c) of
Lemma~\ref{rev:app-kernel-lemma} with
\begin{equation}
r_h(v)=\frac{q(t+h,v)-q(t,v)}{h},
\qquad g(v)=\dot q_t(v).
\label{rev:app-time-difference-quotient}
\end{equation}
These are explicit sufficient conditions for the limit-interchange clause in
Proposition~\ref{rev:prop-generator-sign}.

For a regular conditional kernel version, strict invertibility of $q_t$ gives,
for $\Leb$-almost every $u$,
\begin{equation}
P^X_{t,t+h}\bigl(q_t(u),B\bigr)
=\int_I\one_{B}\bigl(q_{t+h}(v)\bigr)c_{t,t+h}(u,v)\dd v.
\label{rev:app-pushforward-kernel}
\end{equation}
Consequently,
\begin{align}
b_X(t,q_t(u))
&=\lim_{h\downarrow0}
\frac{\int_Iq_{t+h}(v)c_{t,t+h}(u,v)\dd v-q_t(u)}{h}.
\label{rev:app-forward-drift-integral}
\end{align}
Using $\int_Ic_{t,t+h}(u,v)\dd v=1$, the difference quotient in
\eqref{rev:app-forward-drift-integral} is exactly
\begin{align}
&\int_I
\frac{q_{t+h}(v)-q_t(v)}{h}c_{t,t+h}(u,v)\dd v
\label{rev:app-drift-time-term}\\
&\qquad+
\frac{\int_Iq_t(v)c_{t,t+h}(u,v)\dd v-q_t(u)}{h}.
\label{rev:app-drift-state-term}
\end{align}
Lemma~\ref{rev:app-kernel-lemma} and
\eqref{rev:app-time-difference-quotient} show that
\eqref{rev:app-drift-time-term} tends to $\dot q_t(u)$.  Definition
\eqref{rev:app-generator-integral} shows that
\eqref{rev:app-drift-state-term} tends to $\mathcal A_tq_t(u)$.  Hence
\begin{equation}
b_X(t,Q_t(u))
=\partial_tq(t,u)+\mathcal A_tq(t,\cdot)(u)
=(\partial_t+\mathcal A_t)q(t,u).
\label{rev:app-full-drift-identity}
\end{equation}
Because $U_t$ is uniform, the statement is naturally a
$\Leb$-almost-everywhere statement in $u$.  A continuous kernel version extends
it to every interior state where the limits exist.  Strict monotonicity of
$q_t$ then transfers the sign of the right-hand side to the physical state
$x=q_t(u)$, which proves the sign assertion in
Proposition~\ref{rev:prop-generator-sign} after the limiting-mean condition is
imposed separately.
\end{proof}

\subsection{Conditional expectations and observable transition densities}
\label{rev:app-conditional-kernel-proof}

\begin{proof}[Detailed proof of Theorems~\ref{rev:thm-exact-density} and
\ref{rev:thm-exact-mr}, and Proposition~\ref{rev:prop-observable-density}]
For $s<t$, define
\[
H_{s,t}(u)=\int_Iq_t(v)c_{s,t}(u,v)\dd v.
\]
Column normalization and Tonelli's theorem give
\begin{align*}
\int_I|H_{s,t}(u)|\dd u
&\le\int_I\int_I|q_t(v)|c_{s,t}(u,v)\dd v\dd u\\
&=\int_I|q_t(v)|
   \left\{\int_Ic_{s,t}(u,v)\dd u\right\}\dd v\\
&=\int_I|q_t(v)|\dd v<\infty.
\end{align*}
For a bounded $\mathcal F_s$-measurable random variable $Z$, the Markov
property, first for bounded truncations of $q_t$ and then by $L^1$ convergence,
gives
\begin{align*}
\E[ZX_t]
&=\E[Zq_t(U_t)]\\
&=\E\!\left[Z\int_Iq_t(v)c_{s,t}(U_s,v)\dd v\right]
=\E[ZH_{s,t}(U_s)].
\end{align*}
This is the defining test-function identity for
\[
\E[X_t\mid\mathcal F_s]=H_{s,t}(U_s).
\]
Since $X_s=q_s(U_s)$ and $U_s\sim\Leb$, equality
$H_{s,t}(U_s)=q_s(U_s)$ almost surely is equivalent to
$H_{s,t}(u)=q_s(u)$ for $\Leb$-almost every $u$.  This proves both directions
of Theorem~\ref{rev:thm-exact-density} for the fixed pair $(s,t)$; no common
null set over the uncountable collection of time pairs is asserted.  Replacing
the last equality by
\[
H_{s,t}(u)-m=\mathcal R_{s,t}\{q_s(u)-m\}
\]
proves both directions of Theorem~\ref{rev:thm-exact-mr}.

Under the density assumptions of
Proposition~\ref{rev:prop-observable-density}, the substitution
$v=F_t(y)$, $\dd v=p_t(y)\dd y$, in the conditional kernel gives
\[
p^X_{s,t}(x,y)
=c_{s,t}\{F_s(x),F_t(y)\}p_t(y).
\]
The three required identities follow by explicit changes of variables:
\begin{align*}
\int_{E_t^\circ}p^X_{s,t}(x,y)\dd y
&=\int_Ic_{s,t}\{F_s(x),v\}\dd v=1,\\
\int_{E_s^\circ}p_s(x)p^X_{s,t}(x,y)\dd x
&=p_t(y)\int_Ic_{s,t}\{u,F_t(y)\}\dd u=p_t(y),\\
\int_{E_r^\circ}p^X_{s,r}(x,z)p^X_{r,t}(z,y)\dd z
&=p_t(y)\int_Ic_{s,r}\{F_s(x),w\}
                   c_{r,t}\{w,F_t(y)\}\dd w\\
&=p_t(y)c_{s,t}\{F_s(x),F_t(y)\}
=p^X_{s,t}(x,y).
\end{align*}
The first line holds for every selected $x$, the second for Lebesgue-a.e. $y$,
and the third for every selected $x$ and Lebesgue-a.e. $y$.  The final equality
is precisely the density Chapman--Kolmogorov identity; integrating it over an
arbitrary target Borel set yields the corresponding pointwise kernel identity.
\end{proof}

\subsection{The generator martingale theorem and rescaling PDE}
\label{rev:app-generator-martingale-proof}

\begin{proof}[Detailed proof of Theorem~\ref{rev:thm-generator-martingale}]
Write
\[
B(t,u)=(\partial_t+\mathcal A_t)q(t,u).
\]
Membership in the time--space extended generator domain means that, after a
localizing sequence if necessary,
\begin{equation}
L_t=q(t,U_t)-q(0,U_0)-\int_0^tB(s,U_s)\dd s
\label{rev:app-dynkin-local-martingale}
\end{equation}
is a local martingale and the finite-variation integral is locally finite.
If $B=0$ for $\dd t\otimes\Leb(\dd u)$-almost every $(t,u)$, uniformity of
$U_s$ and Tonelli's theorem yield, for every $T<\infty$,
\begin{align*}
\E\int_0^T\one_{\{B(s,U_s)\ne0\}}\dd s
&=\int_0^T\int_I\one_{\{B(s,u)\ne0\}}\dd u\dd s=0.
\end{align*}
Thus the integral in \eqref{rev:app-dynkin-local-martingale} vanishes and
$q(t,U_t)$ is a local martingale.  Class--$(D)$ on each finite horizon makes
the stopped family uniformly integrable, so the local martingale is a true
martingale.

Conversely, suppose $q(t,U_t)$ is a true martingale.  It is also a local
martingale, and subtracting the local martingale $L$ in
\eqref{rev:app-dynkin-local-martingale} shows that
\[
A_t:=\int_0^tB(s,U_s)\dd s
\]
is a local martingale.  It is an absolutely continuous finite-variation
process starting from zero.  A finite-variation local martingale is
indistinguishable from a constant, so $A_t=0$ for all $t$ almost surely and
$B(s,U_s)=0$ for $\dd s\otimes\Pp$-almost every $(s,\omega)$.  Uniformity of
$U_s$ and Tonelli's theorem now give
\[
0=\int_0^T\Pp\{B(s,U_s)\ne0\}\dd s
=\int_0^T\int_I\one_{\{B(s,u)\ne0\}}\dd u\dd s.
\]
Hence $B=0$ for $\dd t\otimes\Leb(\dd u)$-almost every $(t,u)$, proving the
converse in Theorem~\ref{rev:thm-generator-martingale}.
\end{proof}

\begin{proof}[Detailed generator proof of
Proposition~\ref{rev:prop-martingale-rescaling}]
For completeness, let
\[
R_t=\mathcal R_{0,t}=\exp\!\left\{\int_0^t\gamma(a)\dd a\right\}
\]
and define
\[
\widetilde q(t,u)=\frac{q(t,u)-m}{R_t}.
\]
Since $R_t$ is deterministic, $R_t'=\gamma(t)R_t$, and
$\mathcal A_t$ acts only on the state variable,
\begin{align}
(\partial_t+\mathcal A_t)\widetilde q(t,u)
=\frac{1}{R_t}
\left[(\partial_t+\mathcal A_t)q(t,u)
-\gamma(t)\{q(t,u)-m\}\right].
\label{rev:app-rescaled-generator}
\end{align}
Thus the rescaled process is a martingale, subject to the preceding domain and
class--$(D)$ qualifications, exactly when
\eqref{rev:strong-mr-generator-PDE} holds.  At the finite-horizon level,
multiplicativity gives $R_t=R_s\mathcal R_{s,t}$ and therefore
\begin{align*}
\E\!\left[\frac{X_t-m}{R_t}\middle|\mathcal F_s\right]
&=\frac{\mathcal R_{s,t}(X_s-m)}{R_s\mathcal R_{s,t}}
=\frac{X_s-m}{R_s}.
\end{align*}
The reverse calculation proves the converse part of
Proposition~\ref{rev:prop-martingale-rescaling} without appealing only to formal PDE
algebra.
\end{proof}

\subsection{Diffusion transforms and reset mixtures}
\label{rev:app-transform-reset-proofs}

\begin{proof}[Detailed proof of Theorem~\ref{rev:thm-diffusion-transform}]
For Theorem~\ref{rev:thm-diffusion-transform}, It\^o's formula gives
\begin{align*}
\dd\{f(t,R_t)-m\}
&=(\partial_t+\mathcal L)f(t,R_t)\dd t
  +\sigma(R_t)f_x(t,R_t)\dd W_t\\
&=\gamma(t)\{f(t,R_t)-m\}\dd t
  +\sigma(R_t)f_x(t,R_t)\dd W_t.
\end{align*}
With $R_t^\gamma=\exp\{\int_0^t\gamma(a)\dd a\}$, the deterministic product
rule yields
\begin{equation}
\dd\!\left[\frac{f(t,R_t)-m}{R_t^\gamma}\right]
=\frac{\sigma(R_t)f_x(t,R_t)}{R_t^\gamma}\dd W_t.
\label{rev:app-transform-martingale}
\end{equation}
Condition \eqref{rev:transform-integrability} makes the stochastic integral in
\eqref{rev:app-transform-martingale} square integrable on every finite
horizon.  Taking the conditional expectation between $s$ and $t$ and using
$R_t^\gamma/R_s^\gamma=\mathcal R_{s,t}$ gives
\eqref{rev:transform-conditional-mean}.  Strict monotonicity gives the
time-dependent inverse $R_s=f^{-1}(s,Y_s)$ and the transition kernel
\[
P^Y_{s,t}(y,B)
=P^R_{s,t}\bigl(f^{-1}(s,y),f^{-1}(t,B)\bigr),
\]
which proves the Markov claim.  Conversely, the conditional-mean identity
makes the left side of \eqref{rev:app-transform-martingale} a martingale;
write the PDE residual as
\[
 H(t,x)=(\partial_t+\mathcal L)f(t,x)
       -\gamma(t)\{f(t,x)-m\}.
\]
Subtracting the local-martingale part in It\^o's formula shows that
$A_t=\int_0^tH(s,R_s)\dd s$ is both a finite-variation process and a local
martingale.  Hence $A$ is indistinguishable from zero and
$H(t,R_t)=0$ for $\dd t\otimes\Pp$-almost every $(t,\omega)$.  Tonelli's
theorem gives, for each $T<\infty$,
\[
 0=\int_0^T\Pp\{H(t,R_t)\ne0\}\dd t
  =\int_0^T\int_E\one_{\{H(t,x)\ne0\}}
       \Law(R_t)(\dd x)\dd t,
\]
which is precisely \eqref{rev:transform-PDE-law-ae}.  If $H$ is continuous
and $\Law(R_t)$ has full support, $H(t,\cdot)=0$ for almost every $t$; time
continuity then yields the pointwise PDE at all remaining times.  Without
these support and continuity conditions, semimartingale uniqueness alone does
not justify a pointwise PDE on states that the diffusion does not visit.
\end{proof}

\begin{proof}[Detailed proof of Theorems~\ref{rev:thm-mixed-kernel} and
\ref{rev:thm-mixed-mr}]
Finally, write $a=a_{s,r}$ and $b=a_{r,t}$.  Invariance gives
$P_{s,r}\Pi=\Pi P_{r,t}=\Pi$ and $\Pi^2=\Pi$, so the reset-kernel composition
in Theorem~\ref{rev:thm-mixed-kernel} expands as
\begin{align*}
&\{aP_{s,r}+(1-a)\Pi\}
 \{bP_{r,t}+(1-b)\Pi\}\\
&\quad=abP_{s,t}
 +\{a(1-b)+(1-a)b+(1-a)(1-b)\}\Pi\\
&\quad=abP_{s,t}+(1-ab)\Pi
=P^{\mathrm{mix}}_{s,t},
\end{align*}
because $ab=a_{s,t}$.  For the Copula densities the corresponding integral is
\begin{align*}
&\int_I\{ac_{s,r}(u,w)+1-a\}
          \{bc_{r,t}(w,v)+1-b\}\dd w\\
&\quad=ab c_{s,t}(u,v)+a(1-b)+(1-a)b+(1-a)(1-b)\\
&\quad=a_{s,t}c_{s,t}(u,v)+1-a_{s,t},
\end{align*}
where row normalization, column normalization, and the base
Chapman--Kolmogorov identity were used term by term.  This density equality is
for every selected $u$ and a.e. $v$; its integrated kernel equality is for every
$u$ and every Borel target set.  Moreover,
$a_{t,t+h}=1-\lambda(t)h+o(h)$ and
$P_{t,t+h}h_0=h_0+h\mathcal G_th_0+o(h)$ on the generator domain.  Therefore
\begin{align*}
\frac{P^{\mathrm{mix}}_{t,t+h}h_0-h_0}{h}
&=\mathcal G_th_0
 +\lambda(t)\{\Pi h_0-h_0\}+o(1),
\end{align*}
which proves \eqref{rev:mixed-generator}.  The conditional centered first
moment is, directly in integral form,
\begin{align*}
\int_E(y-m)P^{\mathrm{mix}}_{s,t}(x,\dd y)
&=a_{s,t}\int_E(y-m)P_{s,t}(x,\dd y)\\
&\quad+(1-a_{s,t})\int_E(y-m)\pi(\dd y)\\
&=a_{s,t}\mathcal R_{s,t}(x-m),
\end{align*}
and its right derivative at $t=s$ gives
$b_{\mathrm{mix}}(s,x)=b_0(s,x)+\lambda(s)(m-x)$.
\end{proof}
